\documentclass[11pt]{article}


\setlength{\parindent}{0pt}
\usepackage{xltxtra}
\usepackage{hyperref}
\hypersetup{hidelinks}
\usepackage{url}
\urlstyle{tt}
\usepackage{xcolor}
\definecolor{CVBlue}{RGB}{23,110,191}
\usepackage{calc}
\usepackage{graphicx}
\usepackage{tikz}
\usetikzlibrary{calc}
\usepackage{fontspec}
\usepackage{xeCJK}
\usepackage{enumitem}
\CJKsetecglue{} %% 取消中文与数字之间的间隙


%% 主文档字体设置
\setmainfont[
    Path = fonts/Main/,
    Extension = .otf,
    BoldFont = texgyretermes-bold.otf, % 加粗字体
]{texgyretermes-regular.otf} % 正文字体

% 中文字体设置
\setCJKmainfont[
    Path = fonts/hansans/,
    Extension = .ttf,
    BoldFont = NotoSansSC-Bold.ttf, % 加粗字体
]{NotoSansSC-Regular.otf} % 正文字体


\usepackage{relsize}
\usepackage{xspace}

% 使用 fontawesome（部分图标）
\usepackage{fontawesome} 

% A4纸，上下左右边距
\usepackage[
    a4paper,
    left=1.2cm,
    right=1.2cm,
    top=1cm,
    bottom=1cm,
    nohead
]{geometry}

\renewcommand{\baselinestretch}{1.5} % 行间距设为1.5

\usepackage{titlesec}
\usepackage{enumitem}
\setlist{noitemsep} % 取消列表项间的额外间距
%\setlist{nosep} % 取消所有垂直间距
\setlist[itemize]{topsep=0.25em, leftmargin=*}
\setlist[enumerate]{topsep=0.25em, leftmargin=*}

% --- 用于控制【不同项目之间】的垂直距离 ---
\newlength{\interProjectSpacing}
\setlength{\interProjectSpacing}{0.9em} % <--- 在此调整项目之间的距离
\newcommand{\projectsep}{\vspace{\interProjectSpacing}}

% --- 用于控制【项目标题】与下方【项目描述】的距离 ---
\newlength{\intraProjectTitleSep}
\setlength{\intraProjectTitleSep}{0.4em} % <--- 在此调整标题和描述的距离
\newcommand{\titlebreak}{\\[\intraProjectTitleSep]}

% --- 用于控制【项目描述】与下方【要点列表】的距离 ---
\newlength{\intraProjectListTopSep}
\setlength{\intraProjectListTopSep}{0.2em} % <--- 在此调整描述和列表的距离

% =======================================================================


\titleformat{\section}         % 定制 \section 命令 
{\large\bfseries\raggedright} % 将 section 标题设置为大号、粗体且左对齐
{}{0em}                      % 可用于为所有 section 添加前缀（如“章节...”）
{}                           % 可用于在标题前插入代码
[{\color{CVBlue}\titlerule}]  % 在标题后插入一条横线
\titlespacing*{\section}{0cm}{*1.6}{*1.2}



\begin{document}
\pagenumbering{gobble}

%%%% 利用tikz来定位学校Logo (已去除个人照片)
\begin{tikzpicture}[remember picture, overlay] 
    \node[anchor = north west] at ($(current page.north west)+(0.5cm,+1.0cm)$) {\includegraphics[height=6cm]{zju.png}};
\end{tikzpicture}%

\centerline{\LARGE\bfseries{桑 \quad 晟}} 

\vspace{0.5em}
% 使用 tabular 环境让个人信息整齐对齐
\begin{center}
    \normalsize
    \begin{tabular}{l@{\hspace{3em}}l}
        电话：182-6820-5288 & 邮箱：3230103447@zju.edu.cn \\
        出生年月：2005.03 & 籍贯：浙江省长兴县 \\
    \end{tabular}
\end{center}
    
\section{\makebox[\widthof{\faGraduationCap}][c]{\color{CVBlue}\faGraduationCap}\ 教育背景}    
\textbf{浙江大学} \hfill 2023.09 -- 至今\\[0.2em] 
计算机科学与技术学院 - 计科2305 \hfill 浙江杭州 
\begin{itemize}[nosep, topsep=\intraProjectListTopSep]
    \item \textbf{学业成绩：}GPA：4.10/4.3 \quad 百分制：90.42
    \item \textbf{学业排名：}25/152
    \item \textbf{语言水平：}CET-6：622
\end{itemize}

\section{\makebox[\widthof{\faUsers}][c]{\color{CVBlue}\faUsers}\ 科研经历}
\textbf{ZJU 3DV 实验室} \quad 科研实习生 \hfill 2025.10 -- 至今
\begin{itemize}[nosep, topsep=\intraProjectListTopSep]
    \item \textbf{基础知识学习}：系统学习 EECS498 ，复现了 NeRF 架构，并针对体渲染过程实现代码级优化与加速。
    \item \textbf{EQA探索}：在具身仿真环境中对主流算法进行跑通测试与优化尝试；探索并初步设计 Long-term EQA Benchmark ，聚焦智能体长时记忆与多事件推理能力的评估方案。
    \item \textbf{实例级角点检测网络}：构建端到端目标位姿估计流水线，利用三维生成模型与 HCCEPose 脚本渲染并构造合成训练数据集；参考 BoxDreamer 架构，设计并实现了基于 ResNet 的单目实例级检测网络，通过预测目标 8 个角点的 Heatmap 实现角点检测。
    \item \textbf{具身智能与机器人学习}：围绕强化学习与人形机器人 Local Motion 开展文献调研与基础学习，了解 PPO、运动模仿与运动跟踪等基本方法，并调研人形机器人全身运动控制与通用运动 Tracker 的相关研究进展。
\end{itemize}

\section{\makebox[\widthof{\faLaptop}][c]{\color{CVBlue}\faLaptop}\ 核心专业课项目}
\begin{itemize}[nosep, topsep=\intraProjectListTopSep, parsep=0.5ex]
    \item \textbf{MiniSQL}：基于 C++ 从零构建关系型数据库，实现了表堆存储引擎、火山模型 SQL 执行器及支持严格两阶段锁与 DFS 死锁检测的并发控制管理器。
    \item \textbf{操作系统核心实验}：基于 RISC-V 64 架构进行内核开发，完整实现了线程调度与上下文切换、虚拟内存管理、用户态 Shell 交互及底层文件系统设计。
    \item \textbf{计算机体系结构核心实验}：设计并实现了支持 RV32I 指令集的流水线 CPU，集成了异常中断处理与 Cache 缓存机制，并运用计分板算法实现了指令的动态调度。
\end{itemize}

\section{\makebox[\widthof{\faCogs}][c]{\color{CVBlue}\faCogs}\ 技术栈}
\begin{itemize}[nosep]
    \item \textbf{编程语言：} C, C++, Python, Verilog, RISC-V
    \item \textbf{操作系统：} Linux, Windows
\end{itemize}

\section{\makebox[\widthof{\faGraduationCap}][c]{\color{CVBlue}\faList}\ 获奖情况}
\begin{itemize}
    \item 2024-2025学年浙江大学二等奖学金 
    \item 2023-2024学年浙江大学三等奖学金 
    
\end{itemize}
    
\section{\makebox[\widthof{\faFlag}][c]{\color{CVBlue}\faFlag}\ 学生工作}
\textbf{竺可桢学院团委宣传部} \quad 干事 \hfill 2023.09 -- 2024.06

\end{document}